% ══════════════════════════════════════════════════
\chapter{Dashboard --- Web Interface}
% ══════════════════════════════════════════════════

\section{Dashboard Ka Overview}

Dashboard ek \textbf{single HTML page} hai jo browser mein khulta hai (\texttt{http://localhost:3000}). Yeh poore system ka visual interface hai --- sab kuch yahan dikhta hai.

\begin{conceptbox}[Dashboard Kyun Zaroori Hai?]
Bina dashboard ke tumhe terminal/console dekhna padta. Dashboard se ek nazar mein pata chal jata hai:
\begin{itemize}[nosep]
    \item Dono houses kitni bijli use kar rahe hain
    \item Relay ON hai ya OFF
    \item AI kya predict kar raha hai
    \item Blockchain pe kitni transactions hui hain
    \item Koi anomaly toh nahi aayi
\end{itemize}
\end{conceptbox}

\section{Design Choices --- Kyun Aise Banaya}

\begin{table}[H]
\centering
\begin{tabularx}{\textwidth}{l X}
\toprule
\textbf{Choice} & \textbf{Kyun} \\
\midrule
\textbf{Dark Theme} & Energy monitoring systems typically dark hote hain (jaise control rooms). Readings zyada clearly dikhti hain dark background pe. \\
\textbf{Single HTML File} & Setup simple rakhne ke liye. Koi build process nahi, koi framework nahi --- sirf ek file hai, server ke saath serve hoti hai. \\
\textbf{No External JS Libraries} & Chart.js ya React use karne se 200KB+ extra load hota. Hamara custom Canvas chart sirf 100 lines mein hai aur perfectly kaam karta hai. Zero dependencies = zero problems. \\
\textbf{Inter Font} & Professional aur modern dikhta hai. Google Fonts se load hota hai. \\
\textbf{Responsive Layout} & Mobile pe bhi properly dikhta hai (1-column layout), desktop pe multi-column. \\
\bottomrule
\end{tabularx}
\caption{Design Decisions}
\end{table}

\section{Dashboard Ke 8 Sections}

\begin{table}[H]
\centering
\begin{tabularx}{\textwidth}{c l X}
\toprule
\textbf{\#} & \textbf{Section} & \textbf{Kya Dikhata Hai Aur Kyun} \\
\midrule
1 & \textbf{Navbar} & 3 status badges: Blockchain (green/red), AI Engine (green/red), Live Feed (green/red). Ek nazar mein pata chal jata hai ke kya chal raha hai kya nahi. \\
2 & \textbf{Live Monitoring} & House A aur B ki readings --- Voltage, Current, Power, Relay status. Har 500ms pe update hota hai. \\
3 & \textbf{AI Prediction} & 3 cards: Predicted power (aglay 1 min), Relay prediction (ON/OFF hona chahiye), Model accuracy (confidence bar ke saath). \\
4 & \textbf{Chart} & Actual vs Predicted power ka line chart. Blue line = actual, purple dashed = predicted. Isse dikhta hai model kitna close hai. \\
5 & \textbf{ML Model Status} & Model ki technical info: MAE, RMSE, MAPE, version, training date, total predictions count. \\
6 & \textbf{Anomaly Alerts} & Agar koi abnormal reading aaye (voltage spike, negative power, etc.) toh yahan warning dikhti hai. \\
7 & \textbf{Blockchain Stats} & Contract address, total transactions, total energy traded, network info. \\
8 & \textbf{Transactions} & Recent transactions ki list --- har ek mein TX hash, power, energy, timestamp dikhta hai. \\
\bottomrule
\end{tabularx}
\caption{Dashboard Ke 8 Sections}
\end{table}

\section{Data Polling --- Dashboard Data Kaise Leta Hai?}

Dashboard continuously server se data \textbf{fetch} karta rehta hai. 2 alag speeds pe:

\begin{table}[H]
\centering
\small
\begin{tabularx}{\textwidth}{l l X}
\toprule
\textbf{Interval} & \textbf{Kya Fetch Hota Hai} & \textbf{Kyun Itni Speed} \\
\midrule
\textbf{Har 500ms} & Live sensor readings (House A, B) & Real-time feel chahiye --- 500ms se lagta hai ke values ``live'' hain \\
\textbf{Har 5s} & AI prediction, ML status, anomalies, blockchain, transactions & Yeh data slowly change hota hai, 5s kaafi hai \\
\bottomrule
\end{tabularx}
\caption{Dashboard Polling Intervals}
\end{table}

\textbf{Kyun 2 alag speeds?} Kyunke live readings har 500ms pe change hoti hain (sensor data), toh fast polling chahiye. Lekin blockchain stats ya AI prediction har second nahi badalta --- 5 seconds pe ek baar kaafi hai. Isse \textbf{network bandwidth bhi bach'ti hai} aur server pe load bhi kam hota hai.

\section{Custom Canvas Chart}

Dashboard mein ek chart hai jo \textbf{Actual power} (blue line) aur \textbf{Predicted power} (purple dashed line) ko overlay karke dikhata hai.

Yeh chart \textbf{pure Canvas 2D API} se bana hai --- koi library nahi. Iska faida:
\begin{itemize}[nosep]
    \item Zero extra file load (Chart.js = 200KB+, hamara = 0 extra)
    \item Exactly wahi design jo chahiye tha (dark theme matching)
    \item Last 50 data points dikhata hai (older automatically hat jaate hain)
\end{itemize}

Har 5 seconds pe jab nayi prediction aati hai, chart mein ek nayi point add hoti hai --- actual aur predicted dono ki.
